% !TeX root = ../../progress-sem1.tex
% 1.x ที่มาและความสำคัญของปัญหา
\section{ที่มาและความสำคัญของปัญหา}
\thaijustify{
การออกแบบประสบการณ์ผู้ใช้และส่วนต่อประสานผู้ใช้ (UX/UI Design) โดยปกติมีลักษณะเป็นกระบวนการที่วนซ้ำและยึดผู้ใช้เป็นศูนย์กลาง ทว่าในทางปฏิบัตินักออกแบบมักถูกจำกัดด้วยเวลา ประสบการณ์ และทรัพยากร กระบวนการออกแบบเริ่มตั้งแต่การรับโจทย์ รวบรวมความต้องการ ทำความเข้าใจผู้ใช้ นิยามปัญหา ระดมความคิด สร้างต้นแบบ และทดสอบ ในขณะเดียวกันยังต้องทำงานร่วมกับผู้มีส่วนได้ส่วนเสียอื่น ในกระบวนการสร้าง wireframe หรือ prototype มักถูกสร้างขึ้นด้วยมือ กระบวนการนี้จึงใช้แรงงานและเวลามาก~\cite{feng2023designing} ทำให้นักออกแบบสำรวจทางเลือกได้ไม่กว้าง
และเก็บข้อเสนอแนะจากผู้ใช้ได้ไม่เพียงพอ~\cite{suh2025storyensemble}
โอกาสที่ข้อเสนอแนะจะถูกนำมาปรับปรุงงานก่อนส่งมอบจึงลดลง สิ่งนี้ก่อให้เกิดช่องว่างระหว่างผลิตภัณฑ์ที่ส่งมอบกับความต้องการที่แท้จริงของผู้ใช้ ซึ่งเป็นผลมาจากการตัดสินใจเชิงออกแบบที่ถูกกำหนดโดยแรงกดดันด้านเวลามากกว่าความต้องการของผู้ใช้จริง}

\thaijustify{
  งานวิจัยที่ผ่านมาชี้ว่า Generative AI สามารถช่วยลดภาระทางด้านเวลาและแรงงานในการสร้าง prototype และลดปัญหาคอขวดของการทดสอบได้ อย่างไรก็ตาม เครื่องมือที่มีอยู่ในปัจจุบันยังคงกระจัดกระจายและส่วนใหญ่มุ่งเน้นเฉพาะงานใดงานหนึ่งเป็นหลัก ยังไม่ถูกนำมาใช้ตลอดทั้ง UX/UI Development Life Cycle \cite{chung2025ureka, lu2022bridging}
}

% \thaijustify{
%   งานวิจัยที่ผ่านมาชี้ให้เห็นว่า Generative AI สามารถช่วยลดช่องว่างดังกล่าวได้ โดยทำให้แต่ละรอบของการวนซ้ำใช้เวลาน้อยลงและส่งผลให้มีการวนซ้ำอย่างสม่ำเสมอในกระบวนการพัฒนาผลิตภัญฑ์เป็นไปได้ในทางปฏิบัติ ผลการศึกษาแสดงให้เห็นว่า Generative AI สามารถสร้าง wireframe ได้โดยอัตโนมัติจากการอธิบายความต้องการทางการออกแบบด้วยภาษาธรรมชาติ \cite{feng2023designing} และยังสามารถสร้าง high-fidelity prototype ได้ทั้งระดับภาพรวมทั้งระบบและระดับ component \cite{yuan2024towards} ซึ่งช่วยลดต้นทุนด้านเวลาและแรงงานของการออกแบบในแต่ละรอบ นอกจากการทำให้การวนซ้ำทดสอบเกิดบ่อยขึ้นแล้ว สภาพแวดล้อมการออกแบบแบบ forward-backward propagation ยังเปิดโอกาสให้ทีมผู้พัฒนาย้อนกลับไปทบทวนและปรับแก้การตัดสินใจก่อนหน้าได้อย่างเป็นพลวัตตลอดหลายรอบการทำงาน แทนที่จะดำเนินไปข้างหน้าเพียงอย่างเดียว \cite{suh2025storyensemble} แนวทางนี้สอดคล้องกับงานวิจัยด้านการใช้งาน (usability) ที่เป็นที่ยอมรับ ซึ่งแสดงให้เห็นว่าการทดสอบบ่อยครั้งด้วยกลุ่มตัวอย่างขนาดเล็กให้ insight ที่นำไปปรับใช้ได้จริงมากกว่าการทดสอบขนาดใหญ่ที่ทำน้อยครั้งกว่า \cite{nielsen2000five} เมื่อพิจารณาโดยรวม งานวิจัยเหล่านี้ชี้ว่า Generative AI สามารถช่วยลดภาระทางด้านเวลาและแรงงานในการสร้าง prototype และปัญหาคอขวดของการทดสอบได้ อย่างไรก็ตาม เครื่องมือที่มีอยู่ในปัจจุบันยังคงกระจัดกระจายและส่วนใหญ่มุ่งเน้นเฉพาะงานใดงานหนึ่งเป็นหลัก ยังไม่ถูกนำมาใช้ตลอดทั้ง UX/UI DLC \cite{chung2025ureka, lu2022bridging}}

\thaijustify{
เพื่อเพิ่มประสิทธิภาพ UX/UI DLC โครงงานนี้ให้ความสำคัญกับกระบวนการ comparative usability testing เป็นพิเศษ ซึ่งการทำการทดสอบซ้ำหลายรอบ นักออกแบบจำเป็นต้องจัดเตรียม two-alternative UI, รวบรวมและจัดระเบียบข้อเสนอแนะจากผู้ใช้, ตลอดจนตีความผลลัพธ์ที่ได้ก่อนเริ่มการวนซ้ำรอบถัดไปด้วยตนเอง เพื่อสนับสนุนกระบวนการดังกล่าว ผู้พัฒนาจึงนำเสนอแพลตฟอร์มที่\mbox{บูรณาการ} Generative AI ที่มีอยู่เข้ากับกระบวนการ comparative usability testing กล่าวคือ นักออกแบบสามารถอธิบายความต้องการทางการออกแบบจากนั้นแพลตฟอร์มจะสร้าง two-alternative UI สำหรับใช้ในการทดสอบและผู้เข้าร่วมทดสอบสามารถให้ข้อเสนอแนะเชิงบริบทด้วยการเลือก component ที่ต้องการให้ข้อเสนอแนะ และข้อเสนอแนะเหล่านี้จะถูกนำไปรวบรวม กลั่นกรอง โดยปัญญาประดิษฐ์แล้วนักออกแบบตรวจสอบและยืนยันเพื่อสร้าง two-alternative UI ใหม่สำหรับการทำ comparative usability testing ในรอบถัดไป}

\thaijustify{
นอกจากนี้ แพลตฟอร์มดังกล่าวยังมุ่งสร้างขั้นตอนการทำงานร่วมกันระหว่างมนุษย์และปัญญาประดิษฐ์อย่างสมดุล (human-in-the-loop) โดยปัญญาประดิษฐ์จะสรุปและวิเคราะห์ข้อเสนอแนะที่ได้รับจากรอบการทดสอบ เพื่อระบุรูปแบบ ข้อมูลเชิงลึก และประเด็นสำคัญ จากนั้นนักออกแบบจะเข้ามามีบทบาทในการพิจารณาผลการวิเคราะห์ที่ปัญญาประดิษฐ์สร้างขึ้น ใช้การตัดสินใจอย่างมืออาชีพและมีวิจารณญาณในการกำหนดทิศทางการออกแบบ \cite{chung2025ureka} เมื่อนักออกแบบตรวจสอบและยืนยันแนวทางแล้ว Generative AI จะดำเนินการตามที่นักออกแบบกำหนดเพื่อสร้าง two-alternative UI สำหรับรอบการทดสอบถัดไป อีกทั้งยังสร้างรายงานการออกแบบ ที่บันทึกว่าได้รับข้อเสนอแนะใดบ้าง ข้อเสนอแนะนั้นได้รับการตีความโดยปัญญาประดิษฐ์และนักออกแบบตัดสินใจอย่างไร เพื่อเป็นเหตุผลประกอบที่ชัดเจนและมีข้อมูลสนับสนุนสำหรับการสื่อสารกับผู้มีส่วนได้ส่วนเสีย โดยสรุป โครงงานนี้มีเป้าหมายเพื่อแสดงให้เห็นว่าแพลตฟอร์มที่เสริมศักยภาพด้วยการใช้ปัญญาประดิษฐ์ร่วมกับมนุษย์ใน UX/UI DLC สามารถลดระยะเวลาและภาระงานที่ต้องทำด้วยมือ เพิ่มความถี่ของวงรอบการทดสอบ และยกระดับคุณภาพโดยรวมของผลิตภัณฑ์ที่ถูกพัฒนา ตลอดจนเป็นเครื่องมือที่ใช้ตรวจสอบย้อนกลับการตัดสินใจตลอดทั้งกระบวนการได้}
